\section{线性递推与特征根}
\label{sec:04-Discrete-Mathematics/02-Recurrences/02}

0, 1, 1, 2, 3, 5, 8, 13, 21, 34, 55——Fibonacci 数列的前 11 项你或许能背出来。第 100 项呢？写个循环递推 100 步也能算出来，不麻烦。但你想知道它的增长率，想知道有没有一个直接公式能跳过中间所有步骤直接给出第 \(n\) 项，公式里最好只出现 \(n\)，不出现前面的项。

这就不只是计算的问题了。你需要找到这个递推规则的``自然模式''——一种能被递推平移后仍然保持自身形状的东西。

\subsection{指数序列是常系数递推的自然模式}

考虑二阶齐次线性递推

\[
a_n=p\,a_{n-1}+q\,a_{n-2}.
\]

你可能会想：既然递推的系数是常数，每一项都由前两项以固定的权重组合而成，也许序列可以写成某种指数的形式。假设 \(a_n=r^n\) 对所有 \(n\) 成立。代入递推：

\[
r^n=p\,r^{n-1}+q\,r^{n-2}.
\]

两边同时除以 \(r^{n-2}\)（暂设 \(r\neq0\)）：

\[
r^2=p\,r+q,\qquad\text{即}\qquad r^2-p\,r-q=0.
\]

如果某个 \(r\) 满足这个二次方程，那么 \(r^n\) 就精确地满足递推——每一项都是前一项的 \(r\) 倍，递推规则对它来说不是约束，而是描述它本来就有的比例关系。这就是``形状不变''的含义：常系数递推让指数序列成为它的本征模式。

两个不同的根 \(r_1,r_2\) 各自给出一个独立解：\(r_1^n\) 和 \(r_2^n\)。因为递推是线性的，任何线性组合也是解：

\[
a_n=A\,r_1^n+B\,r_2^n.
\]

为什么这就穷尽了所有可能？两个待定常数 \(A,B\) 由两个初值唯一确定——你面对的是一个关于 \(A,B\) 的二元一次方程组，系数矩阵的行列式为 \(r_2-r_1\neq0\)（因为根不同），所以总有唯一解。而递推本身也由初值唯一确定序列，所以叠加这两个模式已经覆盖了所有可能的序列。用线性代数的语言说：递推算子作用在指数模式上只乘一个标量——它们是特征向量——不同模式可线性叠加成整个解空间。

\subsection{Fibonacci 的闭式}

Fibonacci 递推 \(F_n=F_{n-1}+F_{n-2}\) 对应 \(p=q=1\)，特征方程 \(r^2-r-1=0\)，根为

\[
\varphi=\frac{1+\sqrt5}{2},\qquad
\psi=\frac{1-\sqrt5}{2}.
\]

通解 \(F_n=A\varphi^n+B\psi^n\)。代入初值 \(F_0=0,F_1=1\)：

\[
\begin{aligned}
A+B&=0,\\
A\varphi+B\psi&=1.
\end{aligned}
\]

由第一式得 \(B=-A\)，代入第二式 \(A(\varphi-\psi)=1\)。\(\varphi-\psi=\sqrt5\)，所以 \(A=1/\sqrt5,\ B=-1/\sqrt5\)。于是

\[
F_n=\frac{\varphi^n-\psi^n}{\sqrt5}.
\]

这是一个精确公式，不是近似。把 \(n=10\) 扔进去算一下：\(\varphi\approx1.618\)，\(\varphi^{10}\approx122.99\)，\(\psi\approx-0.618\)，\(\psi^{10}\approx0.008\)，所以 \(F_{10}\approx(122.99-0.008)/2.236\approx55.0\)——正是 55。

因为 \(|\psi|<1\)，\(\psi^n\) 随 \(n\) 增大迅速消失到可以忽略。长期增长完全由 \(\varphi^n/\sqrt5\) 主导：\(F_n\approx\varphi^n/\sqrt5\)，增长率就是黄金比例 \(\varphi\approx1.618\)。闭式的价值就在这里——你不需要跑 100 步递推就能一眼看到它的增长速度，也能在不写出前 99 项的前提下说出第 100 项大约是几。

上一节那个``不含连续两个 1 的二进制串''的递推 \(a_n=a_{n-1}+a_{n-2}\) 初值不同（\(a_0=1,a_1=2\)），但特征方程完全相同。它的闭式同样由 \(\varphi^n\) 和 \(\psi^n\) 叠加出来，只是 \(A,B\) 换了一组数。这说明特征根掌管的是递推结构的``固有频率''，初值只决定每种模式的权重。

\subsection{重根会产生多项式因子}

如果特征方程的判别式恰好为零，两根重合：\(r_1=r_2=r\)。你只能写出一个指数解 \(r^n\)。但二阶递推需要两个独立解才能覆盖所有初值条件。第二个解从哪来？

答案是 \(n r^n\)。你可以验证：重根意味着 \(p=2r,\ q=-r^2\)，代入递推：

\[
p\,(n-1)r^{n-1}+q\,(n-2)r^{n-2}
=2r\cdot(n-1)r^{n-1}-r^2\cdot(n-2)r^{n-2}
=\bigl(2(n-1)-(n-2)\bigr)r^n
=n r^n.
\]

为什么偏偏是多项式因子？这个现象在\hyperref[sec:03-Linear-Algebra/05-Eigenvalues-and-Eigenvectors/03]{``对角化失败与 Jordan 形''}中有完全的对应：特征值重复但特征向量不足时，Jordan 块会产生形如 \(n^k\lambda^n\) 的项。这里是完全一样的结构——递推算子的特征空间维度不够，需要广义特征向量来补齐。

重数更高时，解的形态是 \(r^n,\ n r^n,\ n^2 r^n,\ \ldots\)，多项式次数递增。

复根的情况也一并处理。若特征方程有共轭复根 \(\rho e^{\pm i\theta}\)（\(\rho\) 是模长，\(\theta\) 是辐角），实序列的通解可以改写为

\[
\rho^n(A\cos n\theta+B\sin n\theta).
\]

模长 \(\rho\) 控制增长或衰减，辐角 \(\theta\) 控制振荡的频率。不用在复数里兜圈子，写成三角函数就回到了实数世界。

\subsection{非齐次项需要一个特解}

当递推右边多出一个与序列无关的驱动项时：

\[
a_n-p\,a_{n-1}-q\,a_{n-2}=g(n),
\]

总解的结构是``齐次通解 + 一个特解''。齐次通解满足``右边为 0''的情况；特解只负责处理右边的驱动项，不必满足任何初始条件——初始条件最后再用来定通解中的常数。

若 \(g(n)\) 是多项式乘指数，先尝试同类形式做特解。但如果这个形式恰好被齐次通解包含了（即它是特征根对应的模式），那就需要乘上足够的 \(n\) 次幂来区分。这个道理和重根处的 \(n r^n\) 是同一个来源。

举个例子：

\[
a_n-2a_{n-1}=1.
\]

齐次部分 \(a_n-2a_{n-1}=0\) 的通解是 \(C\cdot2^n\)。尝试常数特解 \(a_n=A\)，代入 \(A-2A=1\) 得 \(A=-1\)。所以总解 \(a_n=C\cdot2^n-1\)。

再看一个带现实背景的例子。贷款余额：月利率 \(r>0\)，每月计息后固定还款 \(P\)。理想化的余额递推是

\[
B_{n+1}=(1+r)B_n-P.
\]

特解取常数 \(B_*=P/r\)——即利息正好被还款抵消的平衡点。减去特解后，\(B_n-B_*=(1+r)^n(B_0-B_*)\)。这揭示了两个阈值行为：

\begin{itemize}
\tightlist
\item
  如果 \(B_0>P/r\)：还款抵不过利息，余额反而增长。
\item
  如果 \(B_0<P/r\)：余额递减，在某个 \(n\) 后模型会算出负余额——现实中应在还清时停止递推，公式给出的负数只是未经截断的数学外推。
\end{itemize}

浮动利率、手续费和金额取整都会破坏常系数假设。这个递推揭示的是核心机制和临界条件，不是完整的还款协议。知道它的适用范围，比会用公式重要。

特征根方法统一处理常系数线性递推。碰到非线性、变系数或带取整的递推，需要不同的工具——生成函数、主项分析、或者直接上界估计。方法的选择取决于递推的结构，不是看哪个名字更熟。
